% القسم: المناهج
% الفصل: الاستقراء
% المبحث: المقدمة

\documentclass[../../../include/open-logic-section]{subfiles}

\begin{document}

\olfileid{mth}{ind}{int}

\olsection{مقدمة}

الاستقراء تقنية برهان مهمة تستعمل، بصور مختلفة، في جميع مجالات المنطق
وعلم الحاسوب النظري والرياضيات تقريباً. ونحتاج إليه لإثبات كثير من
النتائج في المنطق.

كثيراً ما يقابل الاستقراء بالاستنباط، ويوصف بأنه الاستدلال من الخاص إلى
العام. فإذا شاهدنا، مثلاً، زمردات خضراء كثيرة ولم نشاهد شيئاً نسميه
زمرداً ولا يكون أخضر، فقد نستنتج أن جميع الزمردات خضراء. وهذا استدلال
استقرائي، لأنه ينتقل من حالات خاصة كثيرة، ككون هذه الزمردة خضراء وتلك
الزمردة خضراء وهكذا، إلى قضية عامة، هي أن جميع الزمردات خضراء. أما
الاستقراء \emph{الرياضي} فهو أيضاً استدلال ينتهي إلى قضية عامة، لكنه
يختلف اختلافاً شديداً عن هذا «الاستقراء البسيط».

على وجه التقريب الشديد، ينتهي البرهان الاستقرائي في الرياضيات إلى أن
جميع الأشياء الرياضية من نوع معين لها خاصية معينة. وفي أبسط الحالات
تكون الأشياء الرياضية التي يعالجها البرهان الاستقرائي أعداداً طبيعية.
ويستعمل البرهان الاستقرائي عندئذ لإثبات أن لجميع الأعداد الطبيعية خاصية
ما، وذلك بإظهار أن:
\begin{enumerate}
    \item العدد $0$ له الخاصية، و
    \item كلما كان للعدد~$k$ الخاصية، كان للعدد~$k+1$ الخاصية أيضاً.
\end{enumerate}
وكثيراً ما يمكن أيضاً استعمال الاستقراء على الأعداد الطبيعية لإثبات
قضايا عامة عن أشياء رياضية يمكن إسناد أعداد إليها. فلكل مجموعة متناهية،
مثلاً، عدد متناهٍ~$n$ من العناصر؛ فإذا استطعنا أن نثبت بالاستقراء أن كل
عدد~$n$ له الخاصية «جميع المجموعات المتناهية ذات الحجم~$n$ هي
\dots»، نكون قد أثبتنا شيئاً عن جميع المجموعات المتناهية.

ويمكن تعميم الاستقراء أيضاً على الأشياء الرياضية
\emph{المعرّفة استقرائياً}. فعبارات لغة صورية، مثل عبارات منطق الرتبة
الأولى، معرّفة استقرائياً. و\emph{الاستقراء البنيوي} طريقة لإثبات نتائج
عن جميع هذه العبارات. والاستقراء البنيوي، بوجه خاص، مفيد جداً وشائع
الاستعمال في المنطق.

\end{document}
